\section{Physical nowcasts}\label{sec:parta}

Before asking whether the signals predict a price, we ask whether they predict the
physical quantities they measure directly: weekly EU arrivals, weekly US loadings
and terminal outages, at one- and two-week horizons. These targets have high
signal-to-noise ratios and published ground truth, and they establish what the
pipeline can and cannot see before any price enters the analysis. The models are
labelled A1 to A5 in the order they were pre-registered in the project plan (Part~A
of the plan is the physical layer, Part~B the spread); A3, a survival model of
voyage time, was specified but not built, for reasons recorded in the log. Four
models were pre-registered and scored walk-forward against two naive nulls, and
each was designed to address the diagnosed weakness of the previous one.

\paragraph{Targets, nulls and scoring.} Time is a weekly grid of Mondays from
2018. For an as-of date $t$ the forecast windows are $W_k = (t + 7(k-1),\, t + 7k]$
days, and the target is a count of cargoes in the window: the conditional EU
arrival count for A1 (laden legs already at sea at $t$ whose first European zone
entry falls in $W_k$) and the US laden-departure count for A2 and A4. The
\term{persistence} null is the last fully elapsed week's count, and the
\term{climatology} null is the trailing four-week mean; both are computable at $t$.
Models are scored by mean absolute error (MAE) in cargoes, and a model
\term{beats} a null when its MAE is lower. We report each model's excess error over
the climatology null,
\begin{equation}
  \text{excess} = \frac{\mathrm{MAE}_{\text{model}} - \mathrm{MAE}_{\text{clim}}}{\mathrm{MAE}_{\text{clim}}},
  \label{eq:excess}
\end{equation}
so that a positive value is worse than the null. Calibration is scored by the share
of outcomes inside the model's central 80\% interval, which should lie between 70
and 90\%. The pre-registered bar was to beat both nulls on one-week MAE and to land
inside that coverage band. Figure~\ref{fig:loadings} shows the loadings series the
US-side models forecast, with both nulls drawn on it.

\begin{figure*}[tb]
\centering
\includegraphics[width=0.95\textwidth]{fig11_loadings}
\caption{Weekly laden departures from US export terminals on the NOAA feed, with
the two objects the nowcasts compete against: the trailing four-week mean (the
climatology null) and the exponentially weighted moving average at the smoothing
constant A4 fitted by maximum likelihood, $\alpha=\aFourAlphaFig$. The series
peaks at \loadingsWeeklyPeak{} cargoes in a week and averages
\loadingsWeeklyMeanTwentyFive{} per week in 2025. The 2020 trough is the COVID
cargo-cancellation wave and the mid-2022 step down is the Freeport outage.}
\label{fig:loadings}
\end{figure*}

\paragraph{A1: duration climatology.} For each open laden leg $i$ of age $a$ at $t$,
the probability that it produces an observed European arrival in $W_k$ is
estimated directly as a count ratio over the matured population of the trailing
year,
\begin{equation}
  p_i(W_k) = \frac{N_{\mathrm{eu}}\!\left(\text{duration} \in (a+u_0,\, a+u_1]\right)}{N(\text{open at age } a)},
  \label{eq:a1}
\end{equation}
where $(u_0, u_1]$ is the window offset. The ratio factorises into an
age-conditional probability that the leg is Europe-bound at all and the duration
distribution of Europe-bound legs, but because both come from one population the
factors telescope into the single ratio above. The weekly count is the
Poisson-binomial sum over open legs, with mean $\sum_i p_i$ and an exact
probability mass function by convolution. A1 has no fitted parameter.

\paragraph{A2: negative-binomial count model.} Weekly US loadings $y$ are modelled
as
\begin{equation}
  \log \lambda = \beta_0 + \sum_j \beta_j x_j, \qquad y \sim \mathrm{NB}(\lambda, k),
  \label{eq:a2}
\end{equation}
with five features fixed in advance together with their expected signs: last
week's count, the trailing four-week mean, the number of vessels in berth, ballast
arrivals in the past week, and the anchorage queue depth. The first two features
are the two nulls, so the scorecard asks whether the three physical features add
anything to the autoregression. Coefficients are refit every week on an expanding
window with a minimum of 104 weeks.

\paragraph{A4: Kalman local level.} The same target is modelled as a latent weekly
rate observed with noise,
\begin{equation}
\begin{aligned}
  x_t &= x_{t-1} + w_t, & w_t &\sim N(0, q), \\
  y_t &= x_t + v_t,     & v_t &\sim N(0, r),
\end{aligned}
\label{eq:a4}
\end{equation}
with $(q, r)$ by exact maximum likelihood, walk-forward. In steady state the
filter's one-step forecast is an exponentially weighted moving average,
$\hat{x}_t = \alpha y_t + (1-\alpha)\hat{x}_{t-1}$, with $\alpha$ set by the ratio
$q/r$ and an effective window of $(2-\alpha)/\alpha$ weeks. A4 therefore asks
whether an optimally tuned moving average beats the fixed four-week one.

\begin{table*}[tb]
\centering\small
\caption{One-week-ahead nowcasts of weekly US loadings (A2, A4) and conditional EU
arrivals (A1), scored walk-forward from 2018 on the NOAA feed. MAE is in cargoes per
week; the excess column is Equation~\eqref{eq:excess}. All three fail the
pre-registered bar of beating both nulls.}
\label{tab:parta}
\begin{tabular}{lrrrrr}
\toprule
model & MAE & persistence & climatology & vs climatology & 80\% cov. \\
\midrule
  A1 arrival-count climatology & 2.314 & 2.217 & 1.814 & +27.5\% & 73.7\% \\
  A2 negative-binomial GLM & 3.961 & 2.551 & 2.152 & +84.0\% & 75.7\% \\
  A4 Kalman local level & 2.104 & 2.694 & 2.088 & +0.7\% & 72.6\% \\
\bottomrule
\end{tabular}

\end{table*}

\begin{figure}[t]
\centering
\includegraphics[width=\linewidth]{fig12_parta}
\caption{One-week MAE of each model against its two nulls, from
Table~\ref{tab:parta}. No model beats the four-week climatology; A4 ties it.}
\label{fig:parta}
\end{figure}

\paragraph{Results.} Table~\ref{tab:parta} and Figure~\ref{fig:parta} give the
one-week results. A1 is \aOneSkill{} worse than the climatology null; its estimator
is built from legs at least ninety days old inside a one-year window, so it is
stale, and its errors are the classic lag pattern of under-shooting on the way up
and over-shooting on the way down. A2 is \aTwoSkill{} worse because it
over-extrapolates the trend: the log link on an expanding window places current
features several standard deviations above the training mean and the exponential
converts that distance into multiplicative overshoot. A2 also falsified one of its
pre-registered signs: ballast arrivals entered with coefficient \aTwoBallastCoef{}
against a prior of positive. A4, the Kalman local-level model, is \aFourSkill{}
from the null, which we treat as a tie. The tie carries information. A4 fitted its
smoothing constant by maximum likelihood and chose $\alpha=\aFourAlpha$, an
effective window of \aFourWindow{} weeks. For a series that is a local level plus
noise, the exponentially weighted moving average is the optimal linear predictor
\citep{muth1960}. The fitted model landing next to a four-week simple average
indicates that the optimum was found, and not that four methods happened to tie. On
this target and at this horizon, the naive mean is close to optimal, and a fifth
method would not have changed the conclusion. The sequence is not a monotone
convergence: A2 was the worst of the four, not the middle rung.

\paragraph{A5: outage detection.} A5 tested Bayesian online change-point detection
\citep{adams2007} on each US terminal's daily laden-departure count against
\aFiveOutages{} labelled outages over \aFiveTerminalYears{} terminal-years. The
observation model is Poisson with a Gamma prior on the rate, and the detector
maintains a posterior over the run length $r_t$ (days since the last change) by
\begin{equation}
\begin{aligned}
  P(r_t = r+1) &\propto P(r_{t-1} = r)\, \pi(y_t \mid r)\,(1-H), \\
  P(r_t = 0)   &\propto \textstyle\sum_r P(r_{t-1} = r)\, \pi(y_t \mid r)\, H,
\end{aligned}
\label{eq:bocpd}
\end{equation}
where $\pi(y_t \mid r)$ is the run's negative-binomial posterior predictive and
$H$ the hazard; it alarms when the run-length-averaged rate falls below a fixed
fraction of the pre-change rate. The nulls are two silence rules that read the
time since the last departure: N1 alarms after a fixed 14~days, N2 after five times
the terminal's own baseline inter-departure interval. Table~\ref{tab:parta5} and
Figure~\ref{fig:outage} give the scorecard. BOCPD did not outperform the
rate-relative rule, because daily binning discards the exact timing that a point
process carries: about half of normal days are zero-count at a two-day cadence, so
each zero is weak evidence. N2 detects with a median delay of \aFiveNTwoDelay{}
days at \aFiveNTwoFar{} false alarms per terminal-year, with recall
\aFiveNTwoRecall{} (\aFiveNTwoHits); every miss is a slow-baseline terminal whose
five-times rule exceeds the 21-day detection window. It is a usable but limited
monitor, and it is the baseline rather than the proposed model. Its use should be
stated plainly: a twelve-day median delay is slower than public news of a terminal
outage, so N2 is an automated cross-check on the reconstruction rather than an
alert a trader would act on.

\begin{table*}[tb]
\centering\small
\caption{Outage detection against labelled terminal outages. N2 is the pre-registered
null and the usable artefact; its recall is the relevant number.}
\label{tab:parta5}
\begin{tabular}{lrrr}
\toprule
detector & recall & median delay (d) & false alarms / terminal-yr \\
\midrule
BOCPD (pre-registered) & 12\% (2/16) & 17 & 0.13 \\
N1 absolute silence (14 d) & 88\% (14/16) & 14 & 2.33 \\
N2 rate-relative silence & 38\% (6/16) & 12 & 0.19 \\
\bottomrule
\end{tabular}

\end{table*}

\begin{figure}[t]
\centering
\includegraphics[width=\linewidth]{fig14_outage}
\caption{Recall against false alarms for the three detectors; marker size grows
with median delay. N1's recall is near-tautological, since a labelled outage is by
construction a gap of at least fourteen days, and its false-alarm rate is the real
story.}
\label{fig:outage}
\end{figure}

\paragraph{A withdrawn test.} A natural extension asks whether skill rises with
horizon, since cargo at sea today determines arrivals two to three weeks ahead. This
was pre-specified before any longer-horizon data was examined, with the premise,
also written down, that the nulls would degrade with horizon. The test extends the
A1 window index to $k=1,\dots,4$ and scores each $W_k$ against the same two nulls.
Extending to four weeks produced a \hOneSkillWFour{} improvement over the best null
at week four (Table~\ref{tab:h1}). The improvement is an artefact, and
Figure~\ref{fig:h1} shows why. A1's target is conditional on legs already at sea at
$t$, whereas both nulls are the last elapsed \emph{unconditional} weekly count. The
two coincide at one and two weeks, because no US-to-EU voyage completes in under
seven days, so nothing arrives in the first week that had not already departed.
Past the voyage tail the conditional truth falls toward zero (mean
\hOneTruthWOne{} arrivals in week one and \hOneTruthWFour{} in week four) while the
nulls continue to predict the full weekly rate (\hOnePersistWFour{} MAE). A1 was
forecasting a smaller quantity, not forecasting better; a forecast of zero would
have scored almost as well. The pre-specified premise fails, and the protocol
treats that outcome as grounds to abandon the test. In the comparable region, the
excess moves from \hOneSkillWOne{} to \hOneSkillWTwo{}, which is the predicted
direction but remains below the null. A valid horizon test needs an unconditional
target, which requires composing A1 with a departure-process model such as A2; that
composition is listed as future work.

\begin{table*}[tb]
\centering\small
\caption{A1 extended to four weekly windows. The skill column at W3 and W4 is an
artefact of a conditional target scored against unconditional nulls; the harness no
longer reports it as a result.}
\label{tab:h1}
\begin{tabular}{lrrrrrr}
\toprule
window & mean truth & A1 MAE & persistence & climatology & ``skill'' & 80\% cov. \\
\midrule
  W1 & 5.89 & 2.324 & 2.226 & 1.822 & $-$27.5\% & 73.5\% \\
  W2 & 5.49 & 2.121 & 2.052 & 1.773 & $-$19.6\% & 76.0\% \\
  W3 & 1.98 & 1.170 & 4.146 & 4.058 & +71.2\% & 85.9\% \\
  W4 & 0.75 & 0.691 & 5.247 & 5.208 & +86.7\% & 94.2\% \\
\bottomrule
\end{tabular}

\end{table*}

\begin{figure}[t]
\centering
\includegraphics[width=\linewidth]{fig13_h1}
\caption{The withdrawn horizon test. Bars are the mean conditional truth in each
window; lines are the MAE of A1 and of the two nulls. Beyond week two the target
empties while the nulls keep forecasting a full week's arrivals, so the apparent
skill measures the target changing meaning.}
\label{fig:h1}
\end{figure}
